\paragraph{AQ-136-CARTAN-SCOPE: what the recovered 1976 primary article proves.}
The entire primary article has been read from its page images:
Cartan, H. (1976). Th\'eories cohomologiques. \emph{Inventiones
Mathematicae, 35}, 261--271. The official GDZ article PDF is
\url{https://gdz.sub.uni-goettingen.de/download/pdf/PPN356556735_0035/LOG_0030.pdf}.
Its first page is a repository cover; physical pages 2--12
are the eleven printed article pages 261--271.

Section 2, p.~263, imposes vertical exactness with constant
degree-zero cycles and horizontal acyclicity in every form
degree. Theorem 1, p.~264, obtains ordinary cohomology.
In Section 4, p.~266, a morphism means a simplicial
$k$-linear cochain map: compatibility with multiplication
is expressly not required. Theorem 2 on that page gives
a unique such integration map to ordinary cochains when
forms vanish above the simplex dimension. Theorem 3,
p.~267, proves multiplicativity on cohomology under a
flatness hypothesis on cycle modules. Section 6, p.~268,
weakens the horizontal axiom. Section 7, pp.~269--271,
constructs the Grothendieck divided-power simplex theory,
obtains $H^n(Gr^{*,r}(X))=H^n(X;\mathbb Z)$ for $r\ge n$
and zero for $r<n$, and identifies its weight-$(r,s)$
cohomology product with $\binom{r+s}{r}$ times cup product.
This is the standard theory $B$ used in this module,
written after eliminating one barycentric coordinate.
The article supplies neither strict PD common sources
nor a terminal-localization theorem for all theories.
Its linear integration theorem cannot be substituted for
either original assertion at cote 136, physical 121.

\paragraph{AQ-136-MULTIPLICATIVE: the canonical comparison preserves the weighted cup product without flatness.}
For every $C\in\mathcal A_k$ and every $D\ge0$, there is
$W_D$ with the following property. For $w,v\ge W_D$,
$p,q\ge0$ and $p+q\le D$, let
\[
 \kappa_w^p:H^p(C^{w,*}(X))\xrightarrow{\ \sim\ }H^p(X;k)
                                                               \tag{136.M1}
\]
denote the canonical comparison, using $T^{[w]}$ to
identify the coefficient line with $k$. Then for every
simplicial set $X$,
\[
 \kappa_{w+v}^{p+q}(\alpha\beta)
   =\binom{w+v}{w}\,
                 \kappa_w^p(\alpha)\smile\kappa_v^q(\beta).
                                                               \tag{136.M2}
\]
No flatness assumption on the terms or cycle modules of
$C$ is needed. The proof uses the underlying cycle
\emph{spaces}, not a possibly nonflat tensor product of
cycle modules. The statement concerns the cohomology
product, and does not produce a strict comparison map.

\begin{proof}
Choose $W_D\ge1$ so that, in every weight at least $W_D$,
the uniform vertical tests hold through degree $D+2$
and all horizontal rows through degree $D+2$ are acyclic.
The finite maximum of their cutoffs exists. It applies
also to $w+v$. We make the comparison, representability,
and normalization steps explicit.

\emph{Relative comparison and connecting maps.}
For a simplicial inclusion $Y\subset X$, put
\[
 C^{w,*}(X,Y)=\ker(C^{w,*}(X)\longrightarrow C^{w,*}(Y)).
\]
The normalized complex of $k[X]/k[Y]$ is free, with
basis the nondegenerate simplices of $X$ not in $Y$.
The complete derived-section proof of AQ-136-COMPARISON
therefore applies to this quotient exactly as it applies
to $k[X]$. It gives canonical relative comparisons with
$H^*(X,Y;k)$ throughout the chosen finite range.

Restriction to $Y$ is surjective in each of the required
cochain degrees. One direct verification is to extend
a normalized chain map to a horizontally acyclic row,
one relative free generator at a time in increasing
chain degree. After the previously prescribed boundary
is mapped, its image is a cycle. Acyclicity, including
$H_0=0$, makes that cycle a boundary; choosing a
preimage defines the value on the new generator.
This proves extension in all degrees. Thus the ordinary
relative short exact sequences are exact in the finite
cochain range needed here. Their comparison with the
derived relative-section triangle shows that the
canonical maps $\kappa$ commute with connecting
homomorphisms. They are normalized in degree zero by
the actual coefficient map $S^w\to C^{w,0}$.

The algebra multiplication also gives relative external
products
\[
 H^p(C^w(X,Y))\otimes H^q(C^v(X',Y'))
 \longrightarrow H^{p+q}
 \bigl(C^{w+v}(X\times X',\,
                Y\times X'\cup X\times Y')\bigr),
                                                               \tag{136.M3}
\]
by pulling back both factors and multiplying. Here and
below the suppressed star denotes the full cochain
complex. The products vanish on the indicated union
because one of their factors vanishes on each piece.
If $a$ is a cocycle on a subspace and $\widetilde a$
is a cochain extension, then its connecting class is
represented by $d\widetilde a$. For a second cocycle $b$,
\[
 d(\widetilde a\,b)=d\widetilde a\,b,
   \qquad
 d(a\,\widetilde b)=(-1)^p a\,d\widetilde b
                  \quad (da=0,\ |a|=p).
                                                               \tag{136.M4}
\]
These formulas prove compatibility of (136.M3) with
connecting maps in the first and second factor, with
the displayed sign. They require no tensor exactness.
The same formulas hold for the ordinary relative cup
and cross products; see Hatcher (2002), p.~218, in
the proof of Theorem 3.15 and the discussion preceding
Theorem 3.18.

\emph{Cycle spaces represent the cohomology.}
Write $Z^r_w=\ker(d:C_*^{w,r}\to C_*^{w,r+1})$.
In the chosen weight range $Z^0_w=\underline{S^w}$
and, for $0\le r<D+1$, there are short exact
sequences of simplicial modules
\[
 0\longrightarrow Z^r_w\longrightarrow C_*^{w,r}
                   \xrightarrow{d}Z^{r+1}_w\longrightarrow0.
                                                               \tag{136.M5}
\]
Each middle object is horizontally acyclic. The homotopy
long exact sequences, beginning with the constant
$Z^0_w$, identify $Z^r_w$ as an Eilenberg--Mac Lane
space $K(S^w,r)$: its only nonzero homotopy group is
$\pi_r=S^w$. This includes all homotopy degrees, since
each middle row is acyclic in every horizontal degree.
Every simplicial abelian group is Kan. The surjection
in (136.M5) is a fibration with contractible total
space, so a map $X\to Z^{r+1}_w$ lifts to that total
space exactly when it is nullhomotopic. The lifting
assertion follows by lifting a nullhomotopy from its
zero endpoint; conversely any lift is nullhomotopic
because its target is contractible. Therefore
\[
 H^r(C^{w,*}(X))=[X,Z^r_w].
                                                               \tag{136.M6}
\]
For $r=0$ the same statement says that closed
zero-forms are locally constant coefficient functions.
This is also the argument of Cartan's Theorem 1,
p.~264, applied only to the finitely many rows in
(136.M5).

The identification of the homotopy group in (136.M5)
is the repeated connecting identification starting
from $S^w$. It agrees with the canonical derived
comparison: both use the coefficient map in degree
zero and commute with the relative connecting maps
just proved. Equivalently, on an oriented relative
$r$-cube both identifications send the repeated
boundary class of a coefficient to that coefficient.
The Eilenberg--Mac Lane representing isomorphism is
determined by this normalization on $\pi_r$: its
fundamental cohomology class corresponds to the
identity on $H_r(K(S^w,r);\mathbb Z)=S^w$.
This last characterization is Hatcher, Theorem 4.57
and its explanation, pp.~393--394. Thus (136.M6)
uses the particular $\kappa$ in (136.M1), rather
than an unspecified additive automorphism.

\emph{The lowest external pairing is forced by coefficients.}
Assume first $p,q>0$. Multiplication gives a pointed
simplicial map
\[
 Z^p_w\times Z^q_v\longrightarrow Z^{p+q}_{w+v}
\]
which vanishes on the wedge, and hence factors through
$Z^p_w\wedge Z^q_v$. The latter smash product has
zero reduced integral homology below degree $p+q$
and
\[
 \widetilde H_{p+q}(Z^p_w\wedge Z^q_v;\mathbb Z)
                         = S^w\otimes_{\mathbb Z}S^v.
                                                               \tag{136.M7}
\]
Indeed each factor has first nonzero reduced integral
homology in its indicated degree, equal to its first
homotopy group (for degree one this is abelianization,
already an abelian group). The reduced integral
K\"unneth sequence has no Tor contribution in degree
$p+q$, and no tensor or Tor contribution below it.
This uses the usual free integral chains of spaces,
not the terms of $C$. Exact primary locators are
Hatcher, Theorem 4.32, p.~366, Theorem 3B.6,
pp.~275--276, and its reduced smash-product version
on p.~276.

By the integral universal coefficient theorem,
\[
 \widetilde H^{p+q}(Z^p_w\wedge Z^q_v;S^{w+v})
   =\operatorname{Hom}_{\mathbb Z}
                      (S^w\otimes_{\mathbb Z}S^v,S^{w+v}).
                                                               \tag{136.M8}
\]
There is no Ext term, since the preceding reduced
homology group vanishes. This is Hatcher, Theorem 3.2,
p.~195. Representability by the target Eilenberg--Mac
Lane space now says that the entire pairing is
determined by its values on cross products of sphere
classes representing $a\in S^w$ and $b\in S^v$.
This argument never uses flatness of these coefficient
groups over $\mathbb Z$, either.

Compute that value using the relative product of
oriented cubes $(I^p,\partial I^p)$ and
$(I^q,\partial I^q)$, with $I=\Delta[1]$.
Here is an explicit connecting induction, including
the relative pairs. Write $u_j(a)$ for the class
corresponding to $a$ in the oriented relative $j$-cube;
for $j=0$ this is the coefficient on a point. For
$j\ge1$ put
\[
 \begin{split}
 K&=I\times I^{j-1},\\
 A&=I\times\partial I^{j-1},\\
 L&=(\partial I\times I^{j-1})\cup A=\partial I^j.
 \end{split}
                                                               \tag{136.M9a}
\]
For $j=1$ the boundary of $I^0$ is empty. The
relative group in degree $j-1$ on $(L,A)$ is two
copies of the group for $(I^{j-1},\partial I^{j-1})$.
This follows directly from compatible sections on
the two disjoint end faces modulo $A$. In the long
exact sequence of the triple $A\subset L\subset K$,
the connecting map sends
$(0,u_{j-1}(a))$ to $u_j(a)$, where the orientation
of $I$ is from its zeroth to its first endpoint.
The same assertion holds in $C$ because the relative
comparison commutes with this connecting map.

Keep the second cube $(I^q,\partial I^q)$ fixed.
In the first-factor connecting step use the product
triple with
$A\times I^q\cup K\times\partial I^q$ and
$L\times I^q\cup K\times\partial I^q$.
The first identity of (136.M4) says that the product
of $\delta(0,u_{j-1}(a))$ with $u_q(b)$ is the
connecting image of
$(0,u_{j-1}(a)\mathbin{\times}u_q(b))$.
This is exactly the same identity and orientation
for ordinary cross product. In the initial case
$j=0$, multiplication by the coefficient $a$ is a
cochain map whose coefficient map is multiplication
in $S$, so naturality of the comparison gives the
claimed value on $u_q(b)$. Induction on $j$ now
proves that the product on the two relative cubes
is ordinary cross product followed by
\[
 S^w\otimes_{\mathbb Z}S^v\longrightarrow S^{w+v},
 \quad (aT^{[w]})\otimes(bT^{[v]})
       \longmapsto ab\binom{w+v}{w}T^{[w+v]}.
                                                               \tag{136.M9}
\]
Surjectivity of restriction, also for these triples,
is the same relative free-chain extension proved
above. Thus each connecting step used here is an
actual cochain calculation for $C$. This proves the required value
on every simple tensor in (136.M7). Such tensors
generate that group, so (136.M8) proves equality of
the universal pairings. Pulling them back along the
two classifying maps and the diagonal of $X$ gives
(136.M2) for $p,q>0$.

If $p=0$, a class is a locally constant coefficient
function. On each connected component multiplication
by its value $aT^{[w]}$ is a natural cochain map
$C^{v,*}\to C^{w+v,*}$ whose map on the coefficient
line is multiplication by the value in (136.M9).
Naturality of the comparison roof therefore gives
(136.M2) in this case. The argument for $q=0$ is
the same, and for $p=q=0$ the formula is simply
the divided-power multiplication in $S$. This treats
every asserted degree, with the single original
weight cutoff $W_D$.
\end{proof}

\paragraph{AQ-136-PD-OPERATIONS: a remaining natural invariant, with its precise torsion limitation.}
Fix $r\ge1$ and $m\ge1$. For every $C\in\mathcal A_k$,
its divided powers induce, for all sufficiently large
$w$, a natural operation
\[
 G_{C,m,w}^{2r}:H^{2r}(X;k)\longrightarrow H^{2rm}(X;k).
                                                               \tag{136.PD1}
\]
The operation is defined using the canonical weight-$w$
and weight-$mw$ comparisons. Strict morphisms of
admissible theories preserve it in their common weight
range. It satisfies
\[
 m!\,G_{C,m,w}^{2r}(a)
       =\frac{(mw)!}{(w!)^m}\,a^m.
                                                               \tag{136.PD2}
\]
Consequently, if multiplication by $m!$ is injective on
$H^{2rm}(X;k)$, then
\[
 G_{C,m,w}^{2r}(a)
       =\frac{(mw)!}{m!(w!)^m}\,a^m.
                                                               \tag{136.PD3}
\]
In particular this holds over a $\mathbb Q$-algebra.
The displayed integer coefficient counts partitions of
$mw$ labelled elements into $m$ unordered blocks of
size $w$, so it is integral for $w\ge1$.

\begin{proof}
The special differential identity shows that
$\gamma_m$ sends closed even forms of bidegree
$(w,2r)$ to closed forms of bidegree $(mw,2rm)$.
It is a simplicial map
$Z^{2r}_w\to Z^{2rm}_{mw}$, preserving zero; it need
not be a linear map. Choose the weight cutoff so that
the two cycle-space descriptions (136.M6) hold.
Two cocycles which define the same cohomology class
are then homotopic as maps from $X$ to $Z^{2r}_w$.
Composing their simplicial homotopy with $\gamma_m$
gives a homotopy in $Z^{2rm}_{mw}$. Hence their
images define the same cohomology class. This proves
well-definedness without presuming that divided powers
of arbitrary boundaries are always boundaries in an
unrestricted PD differential algebra.
Naturality in $X$ is immediate from sections, and
strict PD morphisms commute with the operation.

The identity $m!\gamma_m(x)=x^m$ holds in every PD
algebra. Iterating the product formula (136.M2) gives
the factor
\[
 \prod_{j=1}^{m-1}\binom{(j+1)w}{w}
                          =\frac{(mw)!}{(w!)^m}
\]
for the $m$th power of a weight-$w$ class. Taking a
single large enough cutoff for these finitely many
products proves (136.PD2). Subtract the right-hand
side of (136.PD3) from the operation. Its multiple
by $m!$ is zero, so injectivity proves (136.PD3).
\end{proof}

The ordinary cup-product possibility therefore gives
no component obstruction under the exact admissibility
tests. Higher divided-power operations are well defined
and preserved, but the equation (136.PD2) alone leaves
an $m!$-torsion ambiguity. No pair of admissible theories
with different operations has been constructed here,
and no assertion that this ambiguity is realized is
made. Nor is equality of these operations claimed
sufficient for a strict common source or for vanishing
of all localization loops.

\paragraph{Additional primary reference.}
Hatcher, A. (2002). \emph{Algebraic topology}. Cambridge
University Press. Corrected author PDF:
\url{https://pi.math.cornell.edu/~hatcher/AT/AT.pdf}.
The exact statements used above are Theorem 3.2,
p.~195; the relative connecting-product calculation,
p.~218; Theorem 3B.6 and its reduced relative form,
pp.~275--276; Theorem 4.32, p.~366; and Theorem 4.57,
pp.~393--394. The source statements were inspected;
their full textbook proofs are cited inputs, not
independently reproved here.
